\documentclass{article}
\usepackage{mathtools}
\usepackage{unicode-math}
\setmainfont{Times New Roman}
\setsansfont{Arial}
\setmonofont{Consolas}
\setmathfont{XITSMath-Regular.otf}[BoldFont=XITSMath-Bold.otf]
\usepackage{xeCJK}
\setCJKmainfont{SimSun}
\setCJKsansfont{SimHei}
\setCJKmonofont{SimSun}
\usepackage[version=4]{mhchem}
\usepackage{physics}
\usepackage{xcolor}
\usepackage{cancel}
\usepackage[active,tightpage]{preview}
\PreviewBorder=1pt

\begin{document}
$\symbf{\alpha}$ 中文测试

文章强调，教育、科技、人才是全面建设社会主义现代化国家的基础性、战略性支撑。建设教育强国、科技强国、人才强国具有内在一致性和相互支撑性。要增强系统观念，坚持教育优先发展、科技自立自强、人才引领驱动，统筹推进教育科技人才体制机制一体改革，实现科教兴国战略、人才强国战略、创新驱动发展战略有效联动，形成推动高质量发展的倍增效应。

　　文章指出，要强化教育对科技和人才的支撑作用。科技创新靠人才，人才培养靠教育。要强化高水平研究型大学国家基础研究主力军和重大科技突破策源地作用，建立科技创新与人才培养相互支撑、带动学科高质量发展的有效机制，从国家战略需求中凝练重大科技问题，持续产出原创性、颠覆性科技创新成果。优化高等教育布局，探索国家拔尖创新人才培养新模式。分类推进高校改革发展，引导高校在不同领域不同赛道发挥优势、办出特色。统筹职业教育、高等教育、继续教育，推进职普融通、产教融合、科教融汇，源源不断培养高素质技术技能人才、大国工匠、能工巧匠。推进素质教育，创新教育方法，努力形成有利于创新人才成长的育人环境。

　　文章指出，要构建支持全面创新体制机制，提升国家创新体系整体效能。科学技术是第一生产力、第一竞争力。要完善党中央对科技工作统一领导的体制，健全新型举国体制，强化国家战略科技力量，优化配置创新资源，力争尽早成为世界主要科学中心和创新高地。优化国家科研机构、高水平研究型大学、科技领军企业定位和布局，强化基础研究领域、交叉前沿领域、重点领域前瞻性、引领性布局。激发各类创新主体活力，瞄准世界科技前沿，在加强基础研究、提高原始创新能力上持续用力，在突破关键核心技术、前沿技术上抓紧攻关。
\end{document}